\documentclass[12pt]{article}
\usepackage{amsmath,amssymb,amsthm}
\usepackage[margin=1in]{geometry}

\newtheorem{theorem}{Theorem}
\newtheorem{lemma}[theorem]{Lemma}
\newtheorem{corollary}[theorem]{Corollary}
\newtheorem{definition}[theorem]{Definition}

\title{Problem 491: Double Pandigital Number Divisible by 11}
\author{}
\date{}

\begin{document}
\maketitle

\section{Problem Statement}
We call a positive integer \emph{double pandigital} if it uses all the digits $0$ to $9$ exactly twice (with no leading zero). How many double pandigital numbers are divisible by~$11$?

\section{Notation}
A double pandigital number is a $20$-digit string $d_1 d_2 \cdots d_{20}$ with $d_1 \neq 0$ in which every digit $0,1,\ldots,9$ appears exactly twice. Positions $1,3,5,\ldots,19$ are \emph{odd}; positions $2,4,\ldots,20$ are \emph{even}. Write $O = \sum_{i\;\text{odd}} d_i$ and $E = \sum_{i\;\text{even}} d_i$.

\section{Mathematical Foundation}

\begin{theorem}[Divisibility rule for $11$]
An integer $N$ with decimal representation $d_1 d_2 \cdots d_n$ satisfies $11 \mid N$ if and only if $\sum_{i=1}^{n}(-1)^{i+1}d_i \equiv 0 \pmod{11}$.
\end{theorem}

\begin{proof}
Since $10 \equiv -1\pmod{11}$, we have $10^k \equiv (-1)^k\pmod{11}$. Thus
\[
N = \sum_{i=1}^{n} d_i\cdot 10^{n-i} \equiv \sum_{i=1}^{n} d_i(-1)^{n-i} \pmod{11}.
\]
For $n=20$ (even), $(-1)^{n-i} = (-1)^{-i} = (-1)^i$, so $N \equiv E - O \pmod{11}$. Hence $11 \mid N$ iff $O \equiv E\pmod{11}$.
\end{proof}

\begin{lemma}[Reduction to a single congruence]
The divisibility condition $O \equiv E \pmod{11}$ is equivalent to $O \equiv 1 \pmod{11}$.
\end{lemma}

\begin{proof}
The total digit sum is $O + E = 2(0+1+\cdots+9) = 90$, so $E = 90 - O$. The condition $O \equiv 90 - O\pmod{11}$ becomes $2O \equiv 90 \equiv 2\pmod{11}$. Since $\gcd(2,11)=1$, dividing gives $O \equiv 1 \pmod{11}$.
\end{proof}

\begin{definition}[Assignment vector]
An \emph{assignment vector} is $\mathbf{c} = (c_0,\ldots,c_9) \in \{0,1,2\}^{10}$ with $\sum c_d = 10$, where $c_d$ is the number of copies of digit $d$ placed at odd positions.
\end{definition}

\begin{lemma}[Feasible values]
The odd-position sum $O = \sum d\cdot c_d$ satisfies $0 \le O \le 90$. The values with $O\equiv 1\pmod{11}$ in this range are exactly
\[
O \in \{1,\,12,\,23,\,34,\,45,\,56,\,67,\,78,\,89\}.
\]
\end{lemma}

\begin{proof}
The minimum is $O=0$ (all large digits at even positions) and the maximum is $O=90$. Neither satisfies $O\equiv 1\pmod{11}$. The nine listed values are the complete intersection of the residue class with $[0,90]$.
\end{proof}

\begin{theorem}[Counting formula]
For each valid assignment vector $\mathbf{c}$, the number of leading-zero-free double pandigital numbers is:
\[
N(\mathbf{c}) = \frac{10!}{\prod_d c_d!}\cdot\frac{10!}{\prod_d(2-c_d)!} \;-\; \mathbf{1}_{c_0\ge 1}\cdot\frac{9!}{(c_0-1)!\prod_{d\ge 1}c_d!}\cdot\frac{10!}{\prod_d(2-c_d)!}.
\]
\end{theorem}

\begin{proof}
The first product counts all arrangements of $c_d$ copies of digit~$d$ at odd positions and $2-c_d$ copies at even positions via independent multinomial coefficients. The second term subtracts arrangements with digit~$0$ at position~$1$: fixing one zero there leaves $9$ odd-position slots for the remaining digits, counted by $9!/((c_0-1)!\prod_{d\ge 1}c_d!)$, multiplied by the unchanged even-position count.
\end{proof}

\begin{corollary}
The answer is $\sum_{\mathbf{c}\;\text{valid}} N(\mathbf{c})$.
\end{corollary}

\section{Algorithm}
\begin{enumerate}
\item Enumerate all $(c_0,\ldots,c_9) \in \{0,1,2\}^{10}$ with $\sum c_d = 10$.
\item For each, check $\sum d\cdot c_d \equiv 1 \pmod{11}$.
\item If valid, compute $N(\mathbf{c})$ and accumulate.
\end{enumerate}

\section{Complexity Analysis}
\begin{itemize}
\item \textbf{Time:} $O(3^{10}\cdot 10)$ --- enumerate $59049$ vectors, each in $O(10)$.
\item \textbf{Space:} $O(1)$ auxiliary.
\end{itemize}

\section{Answer}

\[
\boxed{194505988824000}
\]

\end{document}
